% Иллюстрация к выводу статвеса фермионов:
% N фермионов в Z ячейках (Z >= N), не более одного в ячейке (принцип Паули).
% Справа: запрещённая конфигурация с двумя фермионами в одной ячейке.
% Компиляция: pdflatex fermi-cells-figure.tex
\documentclass[tikz,border=8pt]{standalone}
\usepackage[T2A]{fontenc}
\usepackage[utf8]{inputenc}
\usepackage[russian]{babel}
\usetikzlibrary{arrows.meta,decorations.pathreplacing}

\begin{document}
\begin{tikzpicture}[
    note/.style = {gray!30!black},
    ptr/.style  = {gray!70, thin},
  ]

  % ================= ячейки =================
  \draw[semithick] (0,0) -- (10,0);
  \draw[line width=1.2pt] (0,0) -- (0,1.12);
  \draw[line width=1.2pt] (10,0) -- (10,1.12);
  \foreach \x in {1.667,3.333,5,6.667,8.333} {
    \draw[semithick] (\x,0) -- (\x,1.0);
  }

  % ================= фермионы: не более одного в ячейке =================
  % занятость ячеек: 1, 0, 1, 1, 0, 1  (N = 4)
  \foreach \x in {0.83, 4.17, 5.83, 9.17} {
    \shade[ball color=gray!45] (\x,0.23) circle (0.21);
    \draw[gray!60, thin] (\x,0.23) circle (0.21);
  }

  % ================= подписи =================
  \node[note] at (1.2,1.9) {$N$ фермионов};
  \draw[ptr] (1.26,1.64) -- (0.85,0.50);
  \node[note] at (4.3,1.9) {$Z-N$ пустых ячеек};
  \draw[ptr] (3.72,1.66) -- (2.60,0.62);

  \draw[decorate, decoration={brace, mirror, amplitude=6pt}]
        (0,-0.28) -- (10,-0.28) node[midway, below=8pt] {$Z$ ячеек};

  % ================= запрещено: два фермиона в одной ячейке =================
  \draw[semithick] (11.3,0) -- (12.97,0);
  \draw[semithick] (11.3,0) -- (11.3,1.0);
  \draw[semithick] (12.97,0) -- (12.97,1.0);
  \shade[ball color=gray!45] (11.85,0.23) circle (0.21);
  \draw[gray!60, thin] (11.85,0.23) circle (0.21);
  \shade[ball color=gray!45] (12.45,0.23) circle (0.21);
  \draw[gray!60, thin] (12.45,0.23) circle (0.21);
  \draw[red!70!black, line width=1.2pt] (11.35,0.05) -- (12.92,1.05);
  \draw[red!70!black, line width=1.2pt] (11.35,1.05) -- (12.92,0.05);

\end{tikzpicture}
\end{document}
